% !TEX program = pdflatex
% E016 -- evaluation graduee des jurisprudences retournees par G7 -- 2026-08-11
% Statut global : exploration interne. Le pilote train-only ne mesure pas la performance finale de G7.

\documentclass[aspectratio=169,10pt]{beamer}

\usepackage[utf8]{inputenc}
\usepackage[T1]{fontenc}
\usepackage[french]{babel}
\usepackage{lmodern}
\usepackage{booktabs}
\usepackage{array}
\usepackage[table]{xcolor}
\newcolumntype{L}[1]{>{\raggedright\arraybackslash}p{#1}}

\usetheme{Madrid}
\usecolortheme{seahorse}
\setbeamertemplate{navigation symbols}{}
\setbeamertemplate{footline}[frame number]

\definecolor{accent}{RGB}{41,92,155}
\definecolor{accentsoft}{RGB}{230,238,247}
\definecolor{ok}{RGB}{30,130,70}
\definecolor{warning}{RGB}{173,101,0}

\newcommand{\keybox}[1]{%
  \begin{center}\colorbox{accentsoft}{\parbox{0.92\linewidth}{\centering #1}}\end{center}}
\newcommand{\okbox}[1]{%
  \begin{center}\colorbox{green!12}{\parbox{0.92\linewidth}{\centering #1}}\end{center}}
\newcommand{\warnbox}[1]{%
  \begin{center}\colorbox{yellow!18}{\parbox{0.92\linewidth}{\centering #1}}\end{center}}
\newcommand{\status}[1]{\vfill\tiny\textit{Statut scientifique : #1}}

\title[E016 -- JP graduées]{E016 : évaluation graduée des jurisprudences retournées par G7}
\author{Matthieu Kaeppelin}
\institute{Stage FE Recherche}
\date{11 août 2026}

\begin{document}

% 1
\begin{frame}{E016 mesure l'utilité juridique au-delà de l'identité exacte}
  \small
  \begin{columns}[T,onlytextwidth]
    \begin{column}{0.46\textwidth}
      \keybox{\large\textbf{Hit@10 exact}\\[0.35em]
      \normalsize La jurisprudence attendue est-elle présente dans le top--10 ?}
    \end{column}
    \begin{column}{0.46\textwidth}
      \keybox{\large\textbf{Score gradué@10}\\[0.35em]
      \normalsize Les décisions proposées permettent-elles juridiquement de répondre ?}
    \end{column}
  \end{columns}

  \vspace{0.65em}
  \begin{columns}[T,onlytextwidth]
    \begin{column}{0.53\textwidth}
      \textbf{Population évaluée}
      \vspace{0.35em}
      \begin{itemize}\setlength{\itemsep}{0.4em}
        \item \textbf{754 questions} de l'évaluation interne G7 ;
        \item \textbf{$K=10$}, soit \textbf{7 540 positions} ;
        \item \textbf{7 487 couples question--JP uniques} ;
        \item 53 répétitions sur 30 questions : conservées, mais gain nul après la première occurrence.
      \end{itemize}
    \end{column}
    \begin{column}{0.43\textwidth}
      \textbf{Ce que voit le juge}
      \vspace{0.35em}
      \begin{itemize}\setlength{\itemsep}{0.4em}
        \item la question juridique ;
        \item la fiche Step1 déjà calculée pour la décision ;
        \item \textbf{jamais} le rang G7, la JP gold, G8 ni la distance dans le graphe.
      \end{itemize}
    \end{column}
  \end{columns}

  \warnbox{\small E016 complète le Hit@10 ; il ne transforme pas une décision alternative utile en « bonne JP exacte ».}
  \status{E016 exploratoire sur \texttt{eval\_rich\_retrievable\_strict}, déjà consulté ; aucun tuning G7 autorisé.}
\end{frame}

% 2
\begin{frame}{Du jugement aveugle au score audité}
  \scriptsize
  \begin{center}
  \renewcommand{\arraystretch}{1.28}
  \setlength{\tabcolsep}{4pt}
  \begin{tabular}{c L{2.0cm} c L{2.45cm} c L{2.35cm} c L{2.35cm}}
    \colorbox{accentsoft}{\textbf{1. G7}} & top--10 JP par question
    & $\rightarrow$ & \colorbox{accentsoft}{\textbf{2. Juge aveugle}} & $\rightarrow$
    & \colorbox{accentsoft}{\textbf{3. Score fixed-$K$}} & $\rightarrow$
    & \colorbox{green!12}{\textbf{4. Audit avocat}} \\
    & & & A--E ou \texttt{non\_jugeable} & & $(n_A+0{,}5n_B)/10$ & & 100 cas stratifiés et repondérés \\
  \end{tabular}
  \end{center}

  \vspace{0.6em}
  \begin{columns}[T,onlytextwidth]
    \begin{column}{0.54\textwidth}
      \textbf{Grille fermée}
      \vspace{0.35em}
      \begin{tabular}{@{}c L{5.25cm} c@{}}
        \toprule
        \textbf{Classe} & \textbf{Rôle juridique} & \textbf{Gain} \\
        \midrule
        A & répond directement par la règle ou le raisonnement appliqué & 1 \\
        B & apporte une exception, une limite ou une condition utile & 0,5 \\
        C--E & même question sans solution, proximité factuelle/procédurale, ou hors sujet & 0 \\
        NJ & fiche insuffisante pour juger sans inventer & 0 \\
        \bottomrule
      \end{tabular}
    \end{column}
    \begin{column}{0.42\textwidth}
      \okbox{\textbf{Pilote train-only validé}\\[0.3em]
      Job 935516 sur L40S\\
      \textbf{298/298 réponses valides} ; 0 erreur ; 0 JSON invalide}

      \vspace{0.4em}
      \keybox{\textbf{Campagne complète lancée}\\[0.3em]
      Job 935568 sur L40S\\
      \textbf{7 487 couples} à juger}
    \end{column}
  \end{columns}

  \vspace{0.55em}
  \warnbox{\small Le score ne devient interprétable qu'avec la distribution A--E et l'audit avocat ; vigilance particulière sur la frontière B.}
  \status{processus lancé ; aucun résultat final E016 ni validation avocat à ce stade.}
\end{frame}

\end{document}
